\section{Chiral Bar--Cobar Duality and Koszul Duality Beyond Quadratic (Review of Chiral Koszul duality i.e. the Closed Color)}
\label{sec:bar_cobar}

We now develop the bar--cobar duality for $A_\infty$ chiral algebras, extending the quadratic duality theory of Gui--Li--Zeng \cite{GLZ} to the non-quadratic setting and establishing compatibility with our FM calculus framework.

\subsection{Review: GLZ Quadratic Duality for Chiral Algebras}

\subsubsection{Quadratic Chiral Algebras}

\begin{definition}[Quadratic Chiral Algebra]
\label{def:quadratic_chiral}
A \emph{quadratic chiral algebra} is a chiral algebra $\mathcal{A}$ of the form:
\begin{equation}
\mathcal{A} = T_{\text{ch}}(V) / \langle R \rangle,
\end{equation}
where:
\begin{itemize}
\item $T_{\text{ch}}(V)$ is the free chiral tensor algebra on a graded vertex space $V$;
\item $R \subseteq V \otimes_{\text{ch}} V$ is a subspace of \emph{quadratic relations} (degree-2 elements in the chiral tensor product).
\end{itemize}
\end{definition}

\begin{example}[Heisenberg Vertex Algebra]
\label{ex:Heisenberg_quadratic}
The Heisenberg vertex algebra $\mathcal{H}_k$ at level $k$ is generated by $V = \C \cdot J$ with relation:
\begin{equation}
\{J_\lambda J\} = k.
\end{equation}

This is a quadratic chiral algebra with:
\begin{equation}
R = \text{span}\big\{\{J_\lambda J\} - k\big\} \subseteq V \otimes_{\text{ch}} V.
\end{equation}
\end{example}

\subsubsection{GLZ Koszul Dual}

\begin{definition}[Koszul Dual (GLZ)]
\label{def:GLZ_Koszul}
For a quadratic chiral algebra $\mathcal{A} = T_{\text{ch}}(V) / \langle R \rangle$, the \emph{Koszul dual} is:
\begin{equation}
\mathcal{A}^! = T_{\text{ch}}(V^* \otimes \omega_\C[1]) / \langle R^\perp \rangle,
\end{equation}
where:
\begin{itemize}
\item $V^*$ is the graded dual of $V$;
\item $\omega_\C$ is the canonical bundle of $\C$ (shift by 1 in the chiral direction);
\item $R^\perp \subseteq (V^* \otimes \omega_\C[1])^{\otimes 2}$ is the annihilator of $R$ under the natural pairing.
\end{itemize}
\end{definition}

\begin{theorem}[GLZ Maurer--Cartan Control]
\label{thm:GLZ_MC}
(Gui--Li--Zeng \cite{GLZ}, Theorem 4.5) For a Koszul quadratic chiral algebra $\mathcal{A}$, there is a bijection:
\begin{equation}
\Hom_{\text{VA}}(\mathcal{A}, \mathcal{B}) \longleftrightarrow \MC\big(\mathcal{A}^! \widehat{\otimes} \mathcal{B}\big),
\end{equation}
where $\MC(\cdot)$ denotes Maurer--Cartan elements (solutions of $d\alpha + \frac{1}{2}[\alpha, \alpha] = 0$) in the convolution dg Lie algebra.
\end{theorem}

This result parallels Priddy's Koszul duality for associative algebras and Hinich's duality for Lie algebras.

\subsection{Extension to $A_\infty$ Chiral Algebras: The Chiral Bar Construction}

We now extend GLZ duality to $A_\infty$ chiral algebras where $m_{k \geq 3} \neq 0$.

\subsubsection{The Chiral Bar Complex}

\begin{definition}[Chiral Bar Complex with $A_\infty$ Decorations]
\label{def:chiral_bar}
Let $\mathcal{A}$ be an $\OHT$-algebra (Definition \ref{def:ainfty_chiral}), where the $A_\infty$ chiral algebra structure is as in Definition~\ref{def:ainfty_chiral} (Section~\ref{sec:axioms-Ainfty-chiral}). The \emph{chiral bar complex} $\overline{B}_{\text{ch}}(\mathcal{A})$ is:
\begin{equation}
\overline{B}_{\text{ch}}(\mathcal{A}) = \bigoplus_{n \geq 0} \Omega^\bullet_{\log}\big(\FM_{n+1}(\C)\big) \widehat{\otimes} \mathcal{A}^{\otimes (n+1)},
\end{equation}
equipped with differential:
\begin{equation}
d_{\overline{B}} = d_Q + d_{\text{res}} + d_{A_\infty},
\end{equation}
where:
\begin{enumerate}[label=(\roman*)]
\item $d_Q$ is the internal BV-BRST differential acting on $\mathcal{A}$:
\begin{equation}
d_Q(a_0 \otimes \cdots \otimes a_n) = \sum_{i=0}^n (-1)^{|a_0| + \cdots + |a_{i-1}|} (a_0 \otimes \cdots \otimes Qa_i \otimes \cdots \otimes a_n);
\end{equation}

\item $d_{\text{res}}$ is the \emph{boundary residue differential} from FM collision divisors:
\begin{equation}
d_{\text{res}} = \sum_{\substack{S \subseteq \{0, \ldots, n\} \\ |S| \geq 2}} (-1)^{\varepsilon(S)} \Res_{D_S},
\end{equation}
where $\Res_{D_S}$ extracts the residue at boundary divisor $D_S$ (points in $S$ collide);

\item $d_{A_\infty}$ is the \emph{$A_\infty$ grafting differential}, implementing composition via higher operations $m_k$:
\begin{equation}
d_{A_\infty}(a_0 \otimes \cdots \otimes a_n) = \sum_{k \geq 1} \sum_{0 \leq i < i+k \leq n} (-1)^{\varepsilon(i,k)} (a_0 \otimes \cdots \otimes m_k(a_i, \ldots, a_{i+k-1}) \otimes \cdots \otimes a_n).
\end{equation}
\end{enumerate}
\end{definition}

\begin{remark}[Geometric Interpretation]
The three components of $d_{\overline{B}}$ have clear geometric meanings:
\begin{itemize}
\item $d_Q$: Internal differential on observables (quantum equations of motion);
\item $d_{\text{res}}$: Collision/clustering in the holomorphic direction $\C_z$ (chiral OPE);
\item $d_{A_\infty}$: Concatenation in the topological direction $\R_t$ (homotopy coherence).
\end{itemize}
\end{remark}

\subsubsection{Proof of $d_{\overline{B}}^2 = 0$}

\begin{theorem}[$d_{\overline{B}}^2 = 0$]
\label{thm:bar_differential_squared}
The chiral bar differential satisfies $d_{\overline{B}}^2 = 0$.
\end{theorem}

\begin{proof}
We compute:
\begin{equation}
d_{\overline{B}}^2 = (d_Q + d_{\text{res}} + d_{A_\infty})^2 = d_Q^2 + d_{\text{res}}^2 + d_{A_\infty}^2 + \{d_Q, d_{\text{res}}\} + \{d_Q, d_{A_\infty}\} + \{d_{\text{res}}, d_{A_\infty}\}.
\end{equation}

We verify each term vanishes:

\textbf{(i) $d_Q^2 = 0$:} This holds since $Q^2 = 0$ (BV-BRST nilpotency).

\textbf{(ii) $d_{\text{res}}^2 = 0$:} This is the content of the \textbf{Arnold cancellation} (Theorem \ref{thm:Arnold_full_proof}): codimension-2 corners in $\FM_k(\C)$ admit two factorizations (collide $S$ then $T$, or $T$ then $S$), and the residues cancel with appropriate signs:
\begin{equation}
\Res_{D_S} \circ \Res_{D_T} + \Res_{D_T} \circ \Res_{D_S} = 0.
\end{equation}

This is precisely the FM--Stokes--Arnold package developed in Section \ref{sec:FM_calculus} and Appendix \ref{app:FM_Stokes}.

\textbf{(iii) $d_{A_\infty}^2 = 0$:} This is the $A_\infty$ identity (Definition \ref{def:ainfty_chiral}, equation \eqref{eq:ainfty-relation-raw}). Explicitly:
\begin{equation}
\sum_{i+j=k+1} (-1)^{\varepsilon(i,j)} m_i \circ m_j = 0.
\end{equation}

\textbf{(iv) $\{d_Q, d_{\text{res}}\} = 0$:} The BV-BRST differential $Q$ commutes with residues because $Q$ acts fiberwise on $\mathcal{A}$, while $\Res_{D_S}$ acts on the FM base space.

\textbf{(v) $\{d_Q, d_{A_\infty}\} = 0$:} This is the statement that $Q = m_1$ is compatible with higher operations: $Q(m_k(a_1, \ldots, a_k)) = \sum (-1)^{\varepsilon} m_k(a_1, \ldots, Qa_i, \ldots, a_k)$, which follows from the $k=2$ case of the $A_\infty$ identity.
\textbf{(vi) $\{d_{\text{res}}, d_{A_\infty}\} = 0$:} This is the \textbf{cluster factorization identity} (Theorem \ref{thm:cluster_factorization}): the $A_\infty$ composition via $m_k$ is compatible with FM boundary collisions. Explicitly, when points $\{i, i+1, \ldots, i+k-1\}$ collide in $\FM_n(\C)$, the boundary contribution is:
\begin{equation}
\Res_{D_{\{i, \ldots, i+k-1\}}} \sim m_k(a_i, \ldots, a_{i+k-1}),
\end{equation}
which is precisely the action of $d_{A_\infty}$.

Summing all contributions:
\begin{equation}
d_{\overline{B}}^2 = 0 + 0 + 0 + 0 + 0 + 0 = 0.
\end{equation}
\end{proof}

\subsection{The Chiral Cobar Construction}
Dually, we define:

\begin{definition}[Chiral Cobar Complex]
\label{def:chiral_cobar}
For a conilpotent chiral coalgebra $\mathcal{C}$ (encoding distributions on configuration spaces), the \emph{chiral cobar complex} is:
\begin{equation}
\Omega_{\text{ch}}(\mathcal{C}) = \bigoplus_{n \geq 0} \text{Dist}\big(\FM_{n+1}(\C)\big) \widehat{\otimes} \mathcal{C}^{\otimes (n+1)},
\end{equation}
with differential:
\begin{equation}
d_{\Omega} = d_{\text{co}Q} + d_{\text{cores}} + d_{\text{co}-A_\infty},
\end{equation}
where each term is the dual/opposite of the corresponding term in $d_{\overline{B}}$.
\end{definition}

\subsection{Bar--Cobar Adjunction and Twisting Morphisms}

\begin{theorem}[Bar--Cobar Adjunction]
\label{thm:bar_cobar_adjunction}
There is a natural adjunction:
\begin{equation}
\Omega_{\text{ch}} : (\text{conilpotent chiral coalgebras}) \rightleftarrows (\OHT\text{-algebras}) : \overline{B}_{\text{ch}},
\end{equation}
whose unit and counit are defined by universal twisting morphisms satisfying a Maurer--Cartan equation in the convolution dg Lie algebra.
\end{theorem}

\begin{proof}
The proof applies the standard operadic bar--cobar formalism \cite{BV73,MSS02} to the two-colored HT operad $\OHT = C_\ast(W(\SCchtop))$, verifying the additional structure specific to the chiral-topological setting.

\textbf{Step 1: Conilpotence of the bar construction.}
The bar construction $\overline{B}_{\mathrm{ch}}(\mathcal{A})$ is the cofree conilpotent cooperad on $s^{-1}\overline{\mathcal{A}}$ (the desuspension of the augmentation ideal). Conilpotence means: for every element $x \in \overline{B}_{\mathrm{ch}}(\mathcal{A})$, the iterated coproduct $\Delta^{(n)}(x) = 0$ for $n$ sufficiently large. This holds because:
\begin{enumerate}[label=(\alph*)]
\item The bar construction is graded by tree complexity (number of internal edges). The coproduct $\Delta$ decomposes a tree by cutting an internal edge, reducing the complexity by at least 1.
\item The chiral spectral parameters $\lambda_i$ contribute to a secondary grading by pole order. The tameness condition (Definition~\ref{def:tameness}) ensures finite pole order at each stage, preventing infinite iterations.
\end{enumerate}
Together, these give a filtration by total complexity (tree size + pole order) that is bounded below and exhaustive, ensuring conilpotence.

\textbf{Step 2: Convolution Lie algebra and Maurer--Cartan.}
The convolution dg Lie algebra is $\mathfrak{g} = \mathrm{Hom}(\overline{B}_{\mathrm{ch}}(\mathcal{A}), \mathcal{B})$ with:
\begin{itemize}
\item Differential $d_{\mathfrak{g}}(f) = d_{\mathcal{B}} \circ f - (-1)^{|f|} f \circ d_{\overline{B}}$;
\item Lie bracket $[f,g] = f \star g - (-1)^{|f||g|} g \star f$, where the convolution $f \star g$ is defined by $\Delta(x) \mapsto \mu(f \otimes g)(\Delta(x))$, using the cooperad structure $\Delta$ on $\overline{B}_{\mathrm{ch}}$ and the operad composition $\mu$ on $\mathrm{End}_{\mathcal{B}}$.
\end{itemize}
A twisting morphism $\tau: \overline{B}_{\mathrm{ch}}(\mathcal{A}) \to \mathcal{A}^!$ is a Maurer--Cartan element: $d_{\mathfrak{g}}\tau + \frac{1}{2}[\tau,\tau] = 0$. The universal twisting morphism is the projection $\tau_{\mathrm{univ}}: \overline{B}_{\mathrm{ch}}(\mathcal{A}) \twoheadrightarrow s^{-1}\overline{\mathcal{A}}$, and the MC equation for $\tau_{\mathrm{univ}}$ encodes exactly the $A_\infty$ relations \eqref{eq:ainfty-relation-raw}.

\textbf{Step 3: Cluster factorization identity.}
The compatibility $\{d_{\mathrm{res}}, d_{A_\infty}\} = 0$ (where $d_{\mathrm{res}}$ is the residue differential from FM boundary faces and $d_{A_\infty}$ is the $A_\infty$ differential) is the cluster factorization identity (Theorem~\ref{thm:cluster_factorization}): when a cluster of points collides in $\FM_k(\C)$, the resulting boundary term decomposes as a composition of lower-arity operations with spectral substitution $\Lambda_I = \sum_{i \in I} \lambda_i$. This identity is verified by Stokes' theorem on $\FM_k(\C)$ (Section~\ref{sec:FM_calculus}).

\textbf{Step 4: Adjunction.}
The bar and cobar functors form an adjunction by the universal property: giving a map $\Omega_{\mathrm{ch}}(C) \to \mathcal{A}$ (from the free operad on $sC$) is equivalent to giving a twisting morphism $C \to \mathrm{End}_{\mathcal{A}}$, which is equivalent to giving a coalgebra map $C \to \overline{B}_{\mathrm{ch}}(\mathcal{A})$. Naturality in both variables is immediate from the universal constructions.

When $\mathcal{A}$ is quadratic ($m_{k \geq 3} = 0$), the bar complex truncates to trees with binary vertices, recovering the GLZ bar construction~\cite{GLZ}.
\end{proof}

\subsection{GLZ Compatibility: Quadratic Reduction}

\begin{theorem}[GLZ Compatibility]
\label{thm:GLZ_compatibility}
Let $\mathcal{A}$ be a quadratic chiral algebra (Definition \ref{def:quadratic_chiral}) with $m_{k \geq 3} = 0$. Then:
\begin{enumerate}[label=(\roman*)]
\item The chiral bar complex $\overline{B}_{\text{ch}}(\mathcal{A})$ reduces to the GLZ bar construction:
\begin{equation}
\overline{B}_{\text{ch}}(\mathcal{A})|_{m_{k \geq 3} = 0} \cong \overline{B}_{\text{GLZ}}(\mathcal{A});
\end{equation}

\item The Koszul dual $\mathcal{A}^!$ defined via $\Omega_{\text{ch}}(\overline{B}_{\text{ch}}(\mathcal{A}))$ agrees with the GLZ Koszul dual (Definition \ref{def:GLZ_Koszul});

\item The Maurer--Cartan control (Theorem \ref{thm:GLZ_MC}) lifts to our setting:
\begin{equation}
\Hom_{\OHT}(\mathcal{A}, \mathcal{B}) \longleftrightarrow \MC\big(\overline{B}_{\text{ch}}(\mathcal{A}) \widehat{\otimes} \mathcal{B}\big).
\end{equation}
\end{enumerate}
\end{theorem}

\begin{proof}
When $m_{k \geq 3} = 0$, the differential $d_{A_\infty} = 0$, so:
\begin{equation}
d_{\overline{B}} = d_Q + d_{\text{res}}.
\end{equation}

This is precisely the structure of the GLZ bar construction \cite{GLZ}, where:
\begin{itemize}
\item $d_Q$ is the internal differential (Chevalley--Eilenberg differential for Lie algebras, or derived bracket for chiral algebras);
\item $d_{\text{res}}$ encodes the chiral OPE via residues at collision divisors.
\end{itemize}

The GLZ Koszul dual is recovered because, in the quadratic case, the cobar construction $\Omega_{\text{ch}}$ produces exactly the free chiral algebra on $V^* \otimes \omega_\C[1]$ modulo the annihilator relations $R^\perp$.

The Maurer--Cartan bijection follows from the standard argument: morphisms $\mathcal{A} \to \mathcal{B}$ correspond to twisting cochains $\tau : \overline{B}_{\text{ch}}(\mathcal{A}) \to \mathcal{B}$ satisfying $d\tau + \frac{1}{2}[\tau, \tau] = 0$.
\end{proof}

\subsection{Example: Heisenberg Vertex Algebra and Its Koszul Dual}
\begin{example}[Heisenberg and Its Koszul Dual]
\label{ex:Heisenberg_Koszul}
Consider the Heisenberg vertex algebra $\mathcal{H}_k$ at level $k$ (Example \ref{ex:Heisenberg_quadratic}), generated by $J$ with relation $\{J_\lambda J\} = k$.

\textbf{Step 1: Identify as quadratic.}

This is a quadratic chiral algebra:
\begin{equation}
\mathcal{H}_k = T_{\text{ch}}(\C \cdot J) / \langle \{J_\lambda J\} - k \rangle.
\end{equation}

\textbf{Step 2: Compute the Koszul dual.}

The Koszul dual (GLZ, Definition \ref{def:GLZ_Koszul}) is:
\begin{equation}
\mathcal{H}_k^! = T_{\text{ch}}(\C \cdot J^* \otimes \omega_\C[1]) / \langle R^\perp \rangle.
\end{equation}

The annihilator relation $R^\perp$ is computed from the pairing:
\begin{equation}
\langle \{J_\lambda J\}, J^* \otimes J^* \rangle = k.
\end{equation}

Therefore:
\begin{equation}
\mathcal{H}_k^! = CE(\text{Heisenberg Lie algebra}),
\end{equation}
the Chevalley--Eilenberg complex of the Heisenberg Lie algebra (not the vertex algebra!).

\textbf{Step 3: Distinguish from quadratic projection.}

\textbf{Warning}: The Koszul dual $\mathcal{H}_k^!$ is \textbf{not} the same as the \emph{quadratic projection} (taking the associated graded with respect to the PBW filtration). The quadratic projection would yield:
\begin{equation}
\text{gr}(\mathcal{H}_k) = \text{Sym}(\C \cdot J),
\end{equation}
a symmetric algebra (commutative).

In contrast, the Koszul dual $\mathcal{H}_k^!$ is the CE complex, which is \textbf{non-commutative} and has a nontrivial differential encoding the Lie bracket.

This distinction is crucial: Koszul duality exchanges algebraic structures (commutative $\leftrightarrow$ Lie), while the quadratic projection simply "forgets" higher operations.
\end{example}

\begin{remark}[Curvature and Non-Flat Connections]
For curved $A_\infty$ algebras (where $m_0 \neq 0$, i.e., a curvature term), the bar--cobar formalism extends by including a \emph{curved differential}:
\begin{equation}
d_{\overline{B}} = m_0 + d_Q + d_{\text{res}} + d_{A_\infty},
\end{equation}
where $m_0 \in \mathcal{A}$ is a degree-2 element (curvature). The Maurer--Cartan equation becomes:
\begin{equation}
m_0 + d\tau + \frac{1}{2}[\tau, \tau] + \sum_{k \geq 3} m_k(\tau^{\otimes k}) = 0.
\end{equation}

This controls \emph{curved deformations} of chiral algebras, relevant for non-flat connections in quantum field theory.
\end{remark}
